\subsubsection*{Summary and Transition}

\begin{remark*}[Working vocabulary]
This summary reviews argument forms, validity, counterexamples, valid inference
patterns, invalid forms, replacement rules, and the semantic foundation for
ordinary proof moves.
\end{remark*}

\begin{remark*}[Exposition]
An argument form is valid when truth of the premises under an arbitrary truth
assignment forces truth of the conclusion. The central negative test is a
counterexample assignment: one assignment making all premises true and the
conclusion false. This section justified both sides of the method by proving
standard valid forms and refuting common invalid forms.
\end{remark*}

\begin{remark*}[Exposition]
The valid tools now grounded include modus ponens, modus tollens, hypothetical
syllogism, disjunctive syllogism, addition, simplification, conjunction
introduction, constructive dilemma, and destructive dilemma. The invalid forms
of affirming the consequent and denying the antecedent were separated by
explicit counterexample assignments.
\end{remark*}

\begin{remark*}[Transition]
The next section turns to compactness and finite satisfiability as a preview of
a deeper theme: infinite-looking logical questions can sometimes be controlled
by finite witnesses. The semantic foundation developed here remains the base for
that discussion.
\end{remark*}
